\begin{table}[h!]
\centering
\caption{Composición del gasto por nivel de flexibilidad (\% del PIB)}
\label{tab:flexsemiflex}
\begin{tabular}{@{}lcc@{}}
\toprule
\textbf{Categoría} & \textbf{2023} & \textbf{2024} \\
\midrule
\textbf{Gastos inflexibles}        & \textbf{20.1} & \textbf{20.7} \\
\quad Intereses                    & 3.9  & 4.4  \\
\quad Inversión \textsuperscript{\dag}                    & 1.4  & 1.3  \\
\quad Gastos de personal*          & 2.6  & 2.7  \\
\quad SGP                          & 3.3  & 4.0  \\
\quad Pensiones\textsuperscript{1} & 3.6  & 3.6  \\
\quad Salud\textsuperscript{2}     & 2.2  & 2.4  \\
\quad FEPC                         & 1.7  & 1.2  \\
\quad SENA, ICBF y otros \textsuperscript{3} & 0.3 & 0.3  \\
\quad Universidades                & 0.3  & 0.4  \\
\quad Subsidios Servicios Públicos & 0.3  & 0.2  \\
\quad Cesantías y Salud FOMAG      & 0.1  & 0.1  \\
\quad Deudas Entidades             & 0.3  & 0.1  \\
\quad Sentencias                   & 0.1  & 0.1  \\
\midrule
\textbf{Gastos semiflexibles}      & \textbf{2.8}  & \textbf{2.5}  \\
\quad Inversión                    & 1.2  & 0.7  \\
\quad Adq. Bienes y Servicios      & 0.7  & 0.7  \\
\quad Resto de transferencias & 0.9 & 1.2 \\
\quad FOME y FNG                   & 0.0  & 0.0  \\
\bottomrule
\end{tabular}
\par \footnotesize{$\dag$ Para 2024 se asume el mismo monto de inflexibilidades en la inversión que el estimado por \textcite{CARF2024informecongreso} para el año 2023.}
\par \footnotesize{* Incluye reclasificación por contribuciones a la nómina}
\par \footnotesize{$^1$ Incluye bonos pensionales}
\par \footnotesize{$^2$ Incluye 4.4 de  los 9 puntos del antiguo CREE}
\par \footnotesize{$^3$ Corresponde a 4.6 de los 9 puntos del antiguo CREE}
\par \footnotesize{\textit{Fuente: Elaboración propia con base en \textcite{CARF2024informecongreso} y \textcite{minhacienda_balance_2025}.}}
\end{table}
